\section*{Chapter \ref{ch:pl:dataframes} --- Dataframes}

\begin{solution}{exo:pl:dataframes:1}
Closed and labelled on the left: in the bar $[10\!:\!05, 10\!:\!10)$, labelled 10:05. Closed and labelled on the right: in the bar
$(10\!:\!00, 10\!:\!05]$, also labelled 10:05. Same label, different intervals, five minutes apart.
\end{solution}

\begin{solution}{exo:pl:dataframes:2}
Wide: a row per minute (09:31, 09:32), columns AAA and BBB. A minute in which BBB did not trade has a missing value in the BBB column;
melting back and dropping missing values returns the original four rows. A missing cell means \emph{no observation}, not a price of zero,
and code that fills it must say with what.
\end{solution}

\begin{solution}{exo:pl:dataframes:3}
One partition, \texttt{date=5}, and four of the six columns: \texttt{ts} and \texttt{price}, which the query returns, and \texttt{symbol}
and \texttt{date}, which the filter needs.
\end{solution}

\begin{solution}{exo:pl:dataframes:4}
The right-closed bar $k$ is $(kw, (k+1)w]$. For integers, $kw < t \le (k+1)w$ is $kw \le t-1 < (k+1)w$, so $k = \lfloor (t-1)/w \rfloor$.
Truncation toward zero agrees with the floor for non-negative numbers only: for $t - 1 = -7$ and $w = 5$ the floor is $-2$ and the
truncation $-1$, so the times just below zero are put into the bar just above it, which becomes twice as wide. Negative times appear as soon
as times are measured from a reference such as the session open.
\end{solution}

\begin{solution}{exo:pl:dataframes:5}
The filled series has four more returns, each zero, for the empty bars; the return after each gap is the same in both series, but in the
dropped series it spans ten minutes and is treated as a five-minute return. A 20-bar rolling volatility of the dropped series then mixes
five- and ten-minute returns: it overstates the five-minute volatility around the gaps, and its window covers more than 100 minutes of time.
\end{solution}

\begin{solution}{exo:pl:dataframes:6}
Scans, filters and aggregations can each process one batch and pass it on, holding a bounded state. The chapter's pipeline sorts the
trades and joins them as of the quotes, which needs both sides sorted: those operators hold all of their input, and the streaming engine
falls back to the in-memory engine where an operator is not implemented in streaming.
\end{solution}

\begin{solution}{exo:pl:dataframes:7}
On the month of the tests, symbol S003 (7\,960 trades stamped to the second): the three engines return identical bars under each of the
four conventions, and the total volume is the same under all four --- each trade belongs to exactly one bar; the conventions move volume
between bars and rename the bars, they do not create or lose any.
\end{solution}

\begin{solution}{exo:pl:dataframes:8}
Three flaws. The conventions are the library's defaults and unstated: pandas closes and labels bars on the left for most frequencies but on
the right for month, quarter, year and week ends. The returns are computed on a table sorted by time and not grouped by symbol, so each
return compares the bar of one symbol with the previous row, another symbol's bar. And empty bars are either dropped or missing, so returns
span unequal intervals, and the ranks compare them. Compute returns per symbol, state and fix the bar conventions, and decide what an empty
bar is before ranking.
\end{solution}

\begin{solution}{pb:pl:dataframes:1}
\begin{enumerate}
\item Two; the timestamps are truncated to the second, and one second in 300 is a five-minute boundary, so about 375/300 trades are
expected on edges.
\item Four intervals; 600 shares move (each edge trade leaves one bar and enters its neighbour).
\item Five minutes: the left convention names an interval by its start, the right by its end.
\item Dropped: 74 bars and 73 returns; filled: 78 bars and 77 returns.
\item Either, applied everywhere and recorded with the data: width, closed side, label, the treatment of empty bars and the tie rule. With
left labels, the bar's close is known only at label plus width: joining a bar to other data at its label is a look-ahead.
\item Every intermediate table, in full, at each step.
\item Partition pruning, \omterm{def:pl:columnar-formats:pushdown}{projection pushdown} and \omterm{def:pl:columnar-formats:pushdown}{predicate pushdown}.
\item The engine sees the whole computation before running any of it: it can drop unused columns, move filters to the scan and choose
algorithms.
\item The sorts and the as-of join, which need all of their input in order.
\item DuckDB's integer division truncates toward zero; the bar index needs the floor, which differs for negative numbers.
\item One-minute bars, returns, cross-sectional ranks and each symbol's average rank; as-of joined trades and each symbol's mean effective
spread. The answers are rounded to nine decimals and compared for equality before any time is recorded.
\item pandas 895~MB, Polars eager 769, lazy 654, streaming 614, DuckDB 317.
\item Polars streaming is the fastest (0.82~s), DuckDB uses least memory (317~MB).
\item It reads the whole month into memory, turns string columns into Python objects, and each step materialises a new table.
\item Resource usage is preserved across \texttt{exec} on Linux, so the child's \texttt{getrusage} peak can include the memory of the process
it was started from.
\item \textbf{Named result.} Two bar conventions on the same trades differ in 4 of 74 five-minute volumes (600 shares) and label the same
interval five minutes apart, and filling empty bars turns 73 returns into 77; the same research pipeline over a month needs 895~MB beyond
the imports in pandas, 769, 654 and 614~MB in Polars eager, lazy and streaming, and 317~MB in DuckDB, at 0.8 to 1.3~s, with identical
answers.
\item For a month, any of them. For five years (sixty times the data), a lazy engine over the partitions, DuckDB or Polars streaming, or the
pipeline run per day or month where the computation allows it (returns across days need the previous day's last bar).
\item One shared library of idioms (bars, as-of joins, windows, ranks, pivots) with the conventions as arguments and recorded in the
datasets' metadata: width, closed side, label, empty bars, tie rule, units and time zone.
\item Against an engine's default silently changing a result --- a tie rule, a bucketing, a rounding, the treatment of missing values ---
and against a change of behaviour in a new version of an engine.
\item Stated conventions, one implementation tested across engines, and a measured cost.
\end{enumerate}
\end{solution}

\begin{solution}{iq:pl:dataframes:1}
The width, the closed side, the label, the time zone and session times, the treatment of empty bars, the timestamps' resolution (trades on
edges), the tie order for first and last prices, and which trades are included (condition codes, off-exchange prints). Then compare the two
sets bar by bar on the same trades.
\iqlookfor{Conventions first, then a bar-by-bar comparison.}
\end{solution}

\begin{solution}{iq:pl:dataframes:2}
Eager execution runs each operation at once and materialises every intermediate table; lazy execution records a plan and runs it when a
result is requested, after optimisation (pushdowns, pruning, dropping unused columns) and possibly by streaming. It matters when the data
are large compared with memory or read from many files, and when a pipeline computes much it does not return.
\iqlookfor{Materialisation and optimisation.}
\end{solution}

\begin{solution}{iq:pl:dataframes:3}
Bars per symbol (a partition by date and symbol), returns per symbol, then a window rank partitioned by minute with a stated tie rule, in a
lazy or SQL engine over date partitions, one day at a time, keeping each symbol's last bar of the previous day for the first return. About
$3\,000 \times 390 \times 252 \approx 3\times10^8$ rows: a size for an engine that streams, not for an eager in-memory table.
\iqlookfor{Partitioned, streamed, conventions and ties stated.}
\end{solution}

\begin{solution}{iq:pl:dataframes:4}
Read only the columns and partitions needed (a lazy scan), use compact types (categorical symbols, integer prices and times), process by
day where the computation allows it, and move the pipeline to a lazy or SQL engine that streams; measure the peak before and after.
\iqlookfor{Projection, types, chunking, a lazy engine, a measurement.}
\end{solution}

\begin{solution}{iq:pl:dataframes:5}
A contract per idiom (bars, as-of joins, rolling windows, ranks, pivots) with the conventions as required arguments; one implementation per
engine behind it; equivalence tests on data built to hit boundaries and ties, run in continuous integration for every engine version; the
conventions written into each dataset's metadata; benchmarks of time and memory per engine to guide the choice.
\iqlookfor{Explicit conventions and tested equivalence across engines.}
\end{solution}
